% ============================================================================
% Fichier  : front/01_avant_propos.tex
% Version  : 1.0 -- novembre 2026
% ============================================================================

\chapter*{Avant-propos}
\addcontentsline{toc}{chapter}{Avant-propos}
\markboth{Avant-propos}{Avant-propos}

Presque toutes les organisations camerounaises de quelque taille ont un jour
commandé un schéma directeur. Beaucoup dorment dans un tiroir. Le document
existe, il est relié, il a été payé. Personne ne l'a ouvert depuis sa
signature. Les projets qui se font ne sont pas ceux qu'il annonçait.

Ce cours part de ce constat. Un schéma directeur échoue rarement par sa
technique. Il échoue parce qu'il ne dit pas à quelle stratégie il répond,
parce qu'il n'est pas financé ou parce que personne ne le pilote après sa
validation. Le lecteur apprendra donc à produire un document qui tient sur
trois preuves~: il est aligné sur la stratégie de l'organisation, il est
finançable dans le cadre budgétaire qui s'impose à elle, il est gouverné par
des instances qui se réunissent et qui décident.

Le cours clôt la spécialité Management des systèmes d'information. Il
suppose acquis la gouvernance stratégique des SI (UE MSN-2S3-MSI01),
l'architecture d'entreprise et l'urbanisation (UE MSN-2S3-MSI02), COBIT et
ITIL (UE MSN-2S3-MSI03), la gestion de projet (UE MSN-2S3-MGT01) et l'audit
des SI (UE MSN-2S3-MSI05). Il s'en sert sans les répéter. Il laisse au cours
de conduite du changement (UE MSN-2S4-MSI04) l'adhésion des personnes.

Trois partis pris le traversent. Le premier est de lire les théories dans
leurs textes d'origine, avec leurs dates et leurs limites, parce qu'un modèle
cité de seconde main devient vite un slogan. Le deuxième est de placer le
droit avant la méthode~: une administration n'achète pas, ne finance pas et
ne décide pas comme une entreprise, et un schéma directeur qui l'ignore est
inapplicable dès sa première ligne budgétaire. Le troisième est de coter ce
que l'on sait. Chaque constat du diagnostic et chaque exigence d'audit reçoit
une cote \OUI{}, \NON{} ou \NEK{}, appuyée sur une pièce. Un diagnostic qui
ne distingue pas ce qui est établi de ce qui est supposé fonde une cible sur
du sable.

Avant la méthode, deux parties situent le sujet. La première présente le
contexte international~: les normes et les référentiels, les doctrines
publiques, les stratégies numériques de l'Afrique. La seconde présente le
contexte national~: la stratégie numérique du Cameroun et le cadre juridique
qui s'impose à un schéma directeur. Les textes nationaux sont souvent
l'internalisation de règles venues d'ailleurs, d'où cet ordre.

Les trois organisations qui accompagnent le lecteur, l'ANPP, LiZbiZ et la
BANQUE SANAGA, sont fictives, comme leurs dirigeants, leurs projets, les
dates et les montants. Le cabinet AUDIX CONSEIL l'est aussi. Les théories,
les normes et les textes cités sont exacts. Les normes sous licence sont
paraphrasées.

Le livre est accompagné d'un livret des corrigés et de NDAP SUKULU, où les
fiches de questions et les dix laboratoires du cours permettent de pratiquer
chaque chapitre.

\bigskip
\hfill\begin{minipage}{0.5\textwidth}
\raggedleft
\IfFileExists{../../Logos/signature_avant_propos.png}{\includegraphics[height=1.6cm]{../../Logos/signature_avant_propos.png}}{\textit{MINKA MI NGUIDJOÏ Thierry Emmanuel}}\\
Yaoundé, novembre 2026
\end{minipage}
